\chapter{Worked example: adding a Device Type}\label{ch:device-type}

This chapter is not a description of a Device Type. It is a change made to this codebase, and every
line in it was executed: \texttt{tests/test\_manuals.py} copies \texttt{src/} into a temporary
directory, puts the driver below at \texttt{src/ate/devices/echo.py}, applies the three edits shown
here character for character, and runs \texttt{docs/manuals/examples/developer/echo.tslua} headless
against that copy. If the manual and the code ever part company, that test fails.

A Device Type has three parts, and all three have to exist: the driver that opens it, the registry
line that makes the name mean that driver, and the enum member that lets a schedule write the name
down. Two of them are in the same file.

\section{Step 1: the driver}

A driver is a subclass of :class:`ate.devices.registry.Driver`. It declares the Device Type name,
the fields an \texttt{ATE.open} may pass, and the fields \texttt{ATE.read}/\texttt{ATE.write} may
pass; the registry validates every one of them generically through :func:`validate_params`, which
raises a :class:`ValueError` naming the offending field before anything connects.

\lstinputlisting[language=ATEPython, caption={\texttt{docs/manuals/examples/developer/echo_driver.py}}]{../examples/developer/echo_driver.py}

Three things in that file are not decoration:

\begin{itemize}
  \item \texttt{fields} is the whole of the open-time contract. An undeclared key is refused with
    \texttt{unknown field 'x' for ECHO device}, so a schedule typo never reaches a driver.
  \item \texttt{read\_fields} and \texttt{write\_fields} are \texttt{READ\_FIELDS} and
    \texttt{WRITE\_FIELDS} from \texttt{src/ate/devices/stream.py}. Every Device Type takes the same
    read and write fields; a Device with a different contract would declare its own dict, and the
    ATE wrappers validate against the driver carried on the handle.
  \item \texttt{partial\_read\_is\_truncation = True} on the handle is what makes a cut-short read
    say \texttt{datagram truncated to N byte(s)} instead of \texttt{timed out after N byte(s)}. The
    engine asks the Device which it is (packet 32b review): there is no stream and no deadline here,
    so claiming a timeout would be a lie.
\end{itemize}

\section{Step 2: register the driver}

The registry is built once per run, in \texttt{LuaHost.\_\_init\_\_}
(\texttt{src/ate/engine/lua_host.py}). Add one line beside the others:

\begin{lstlisting}[language=ATEPython]
        self._devices.register("TCP", TcpDriver)
        self._devices.register("SERIAL", SerialDriver)  # packet 18a: pyserial ports, same API
        self._devices.register("UDP", UdpDriver)  # packet 18b: datagrams, one datagram per message
        self._devices.register("VISA", VisaDriver)  # packet 19: pyvisa resources, read on demand (no reader)
        self._devices.register("ECHO", EchoDriver)  # the developer manual's worked example
\end{lstlisting}

and the import above it:

\begin{lstlisting}[language=ATEPython]
from ..devices import DeviceRegistry, EchoDriver, TcpDriver
\end{lstlisting}

Without this line the enum member is a name with nothing behind it: \texttt{ATE.open(DeviceType.ECHO,
\{label = "x"\})} reaches the registry, which raises \texttt{unknown DeviceType 'ECHO'}. That test is
run too --- see the last section of this chapter.

\section{Step 3: the enum member}

The flat enum tables a schedule can see live in \texttt{\_ENUMS} in the same file:

\begin{lstlisting}[language=ATEPython]
    "DeviceType": {
        "TCP": {"TCP": "TCP"},
        "SERIAL": {"SERIAL": "SERIAL"},
        "UDP": {"UDP": "UDP"},
        "VISA": {"VISA": "VISA"},
        "ECHO": {"ECHO": "ECHO"},
    },
\end{lstlisting}

Each entry maps the Lua table name to its members; :func:`_device_type_name` (also in
\texttt{lua_host.py}) turns a member back into the string the registry looks up, and it is the one
that raises \texttt{deviceType must be a DeviceType enum member} when a schedule passes a bare
string or a number.

\section{Step 4: run it}

The schedule below is the one the test runs. It writes a terminator, leaves a Device open at the end
of the run, and cuts a read with \texttt{size} --- three things a driver has to get right, and each
of them shows up in the record.

\lstinputlisting[language=ATELua, caption={\texttt{docs/manuals/examples/developer/echo.tslua}}]{../examples/developer/echo.tslua}

To run it yourself against a copy of the source with the driver in place, put the runner from
chapter~\ref{ch:run-session} in \texttt{runner.py} and start it with the copy of \texttt{src} on
\texttt{PYTHONPATH}:

\begin{lstlisting}[language=ATEShell]
python runner.py /tmp/station/config/config.json "Developer: echo Device"   # PYTHONPATH=<the copy of src>
\end{lstlisting}

\section{What the run proves}

The test asserts the run ends \texttt{PASS} / \texttt{COMPLETED} with all three steps
\texttt{PASS}, and then reads the traces the engine recorded:

\begin{center}
\small
\begin{tabular}{p{5.4cm} p{7.6cm}}
\toprule
\textbf{Line in the Run log} & \textbf{What it proves} \\
\midrule
\texttt{open ECHO manual-example ok (0.0 ms)} & the registry wrapped the new driver like any other,
  and traced it \\
\texttt{write ECHO manual-example 6 bytes "hello\textbackslash n" (0.0 ms)} & the terminator is part
  of the counted payload \\
\texttt{read ECHO sized 4 bytes "ABCD" partial (0.0 ms)} & \texttt{size} cut the read short, and the
  handle reported it as partial \\
\texttt{ATE.read: datagram truncated to 4 byte(s)} & the Device said it was a truncation, not a
  timeout \\
\texttt{device ECHO left-open left open at end of run} (WARNING) & the run-end teardown closed a
  Device the schedule never closed \\
\bottomrule
\end{tabular}
\end{center}

Two more assertions are the negative cases: a Device Type the registry has never registered raises
\texttt{unknown DeviceType 'ECHOX'} from \texttt{DeviceRegistry.open\_device}, and the schedule run
with only the enum member added and no register line fails the run with \texttt{deviceType must be a
DeviceType enum member}. A Device Type is only real when all three parts are there.

\section{What else a new Device Type touches}

\begin{itemize}
  \item A test file named after the packet that added it
    (\texttt{tests/test\_device\_abort\_32b.py} is the style), with the
    \texttt{LoopbackPeer}/\texttt{UdpLoopbackPeer} fixtures in \texttt{tests/\_fixtures.py} if the
    Device needs a peer.
  \item The Device Type list in both test-writer manuals:
    \texttt{missing\_from(MANUALS["test-writer"].texts(), device\_types())} is derived from the
    drivers that exist, so a driver nobody documented is a failing suite, not an omission somebody
    might notice later.
  \item \texttt{docs/manuals/llm/test-writer.md}: a schedule writer reading only that file needs to
    know the Device exists and what its fields accept.
\end{itemize}

\cando{add a Device Type with a driver, a register line and an enum member, and show it in the record.}
